% !PRINTED 199
\weberpage{199}
%<<CONT
\gesperrt{cients $M$ be represented by one of the two
expressions}
\[
\Sigma \, M D^\gamma H^\alpha f^\beta, \quad Q \, \Sigma \, M D^\gamma H^\alpha f^\beta.
\]
The invariants are contained among the covariants as a
special case, and are all powers of $D$.

\begin{center}
\S.\,64.

\gesperrt{Biquadratic Forms.}
\end{center}

We now pass on to the consideration of the following biquadratic form:
\weq{1}{f(x, y) = a_0 x^4 + a_1 x^3 y + a_2 x^2 y^2 + a_3 x y^3 + a_4 y^4}

We have first the \gesperrt{Hessian} determinant as a covariant of the fourth order, which, in order to free it of a numerical
factor, we divide by 3:
\weq{2}{H = \frac{1}{3}
\begin{vmatrix}
12 a_0 x^2 + 6 a_1 x y + 2 a_2 y^2, & 3 a_1 x^2 + 4 a_2 x y + 3 a_3 y^2 \\
3 a_1 x^2 + 4 a_2 x y + 3 a_3 y^2, & 2 a_2 x^2 + 6 a_3 x y + 12 a_4 y^2
\end{vmatrix},}
or
\[
H = A_0 x^4 + A_1 x^3 y + A_2 x^2 y^2 + A_3 x y^3 + A_4 y^4,
\]
in which
\[
\begin{aligned}
A_0 &= 8 a_0 a_2 - 3 a_1^2, & A_4 &= 8 a_2 a_4 - 3 a_3^2, \\
A_1 &= 24 a_0 a_3 - 4 a_1 a_2, & A_3 &= 24 a_1 a_4 - 4 a_2 a_3, \\
A_2 &= 48 a_0 a_4 + 6 a_1 a_3 - 4 a_2^2, & &
\end{aligned}
\]
and a covariant of the sixth degree
\weq{3}{T = \frac{1}{12}
\begin{vmatrix}
f'(x), & f'(y) \\
H'(x), & H'(y)
\end{vmatrix}.}

The invariants of the biquadratic form are most simply
formed, following \S.\,61, from the differences of the roots.

We thus obtain an invariant of the second order and one
of the third order in the coefficients:
\weq{4}{A = \tfrac{1}{2} a_0^2 [(12)^2(34)^2 + (13)^2(24)^2 + (14)^2(23)^2],}
\weq{5}{B = a_0^3 \, \Sigma \, (12)^2(34)^2(13)(42),}
expressions which can be written more perspicuously if one
sets
\weq{6}{U = (12)(34), \quad V = (13)(42), \quad W = (14)(23)}
namely
%CONT>>
